Náhradový poměr (důchod jako~\% čistého příjmu) výrazně klesá s~rostoucí mzdou: nízkopříjmoví zaměstnanci dosahují náhradového poměru přes \SI{60}{\percent}, zatímco vysoko­příjmoví pracovníci jen kolem \SIrange{20}{30}{\percent}.
Tato solidaritní redistribuce je záměrnou součástí důchodové formule a~zajišťuje, že minimální důchod zůstává sociálně přijatelný i~pro pracovníky s~přerušovanou kariérou nebo nízkými mzdami --- skupiny, které jsou v~segmentovaném trhu práce bez silného kolektivního pokrytí nejvíce ohroženy.
Přetrvávající mzdová disper­ze v~ČR (obr.~\ref{fig:rscp_sector_percentiles}) však znamená, že čím větší je podíl nízkopříjmových pracovníků, tím vyšší je implicitní přerozdělovací zátěž systému; bez mzdové konvergence tak narůstá strukturální deficitnost průběžného financování.
Posílení kolektivního vyjednávání na odvětvové úrovni, které by comprimovalo mzdové rozdělení zdola, by tedy zároveň snižovalo redistribuční nároky na důchodový systém a~zvyšovalo jeho dlouhodobou udržitelnost.

Náhradový poměr (důchod jako~\% čistého příjmu) výrazně klesá s~rostoucí mzdou: nízkopříjmoví zaměstnanci dosahují náhradového poměru přes \SI{60}{\percent}, zatímco vysoko­příjmoví pracovníci jen kolem \SIrange{20}{30}{\percent}.
Tato solidaritní redistribuce je záměrnou součástí důchodové formule a~zajišťuje, že minimální důchod zůstává sociálně přijatelný i~pro pracovníky s~přerušovanou kariérou nebo nízkými mzdami --- skupiny, které jsou v~segmentovaném trhu práce bez silného kolektivního pokrytí nejvíce ohroženy.
Přetrvávající mzdová disper­ze v~ČR (obr.~\ref{fig:rscp_sector_percentiles}) však znamená, že čím větší je podíl nízkopříjmových pracovníků, tím vyšší je implicitní přerozdělovací zátěž systému; bez mzdové konvergence tak narůstá strukturální deficitnost průběžného financování.
Posílení kolektivního vyjednávání na odvětvové úrovni, které by comprimovalo mzdové rozdělení zdola, by tedy zároveň snižovalo redistribuční nároky na důchodový systém a~zvyšovalo jeho dlouhodobou udržitelnost.

Náhradový poměr (důchod jako~\% čistého příjmu) výrazně klesá s~rostoucí mzdou: nízkopříjmoví zaměstnanci dosahují náhradového poměru přes \SI{60}{\percent}, zatímco vysoko­příjmoví pracovníci jen kolem \SIrange{20}{30}{\percent}.
Tato solidaritní redistribuce je záměrnou součástí důchodové formule a~zajišťuje, že minimální důchod zůstává sociálně přijatelný i~pro pracovníky s~přerušovanou kariérou nebo nízkými mzdami --- skupiny, které jsou v~segmentovaném trhu práce bez silného kolektivního pokrytí nejvíce ohroženy.
Přetrvávající mzdová disper­ze v~ČR (obr.~\ref{fig:rscp_sector_percentiles}) však znamená, že čím větší je podíl nízkopříjmových pracovníků, tím vyšší je implicitní přerozdělovací zátěž systému; bez mzdové konvergence tak narůstá strukturální deficitnost průběžného financování.
Posílení kolektivního vyjednávání na odvětvové úrovni, které by comprimovalo mzdové rozdělení zdola, by tedy zároveň snižovalo redistribuční nároky na důchodový systém a~zvyšovalo jeho dlouhodobou udržitelnost.

Náhradový poměr (důchod jako~\% čistého příjmu) výrazně klesá s~rostoucí mzdou: nízkopříjmoví zaměstnanci dosahují náhradového poměru přes \SI{60}{\percent}, zatímco vysoko­příjmoví pracovníci jen kolem \SIrange{20}{30}{\percent}.
Tato solidaritní redistribuce je záměrnou součástí důchodové formule a~zajišťuje, že minimální důchod zůstává sociálně přijatelný i~pro pracovníky s~přerušovanou kariérou nebo nízkými mzdami --- skupiny, které jsou v~segmentovaném trhu práce bez silného kolektivního pokrytí nejvíce ohroženy.
Přetrvávající mzdová disper­ze v~ČR (obr.~\ref{fig:rscp_sector_percentiles}) však znamená, že čím větší je podíl nízkopříjmových pracovníků, tím vyšší je implicitní přerozdělovací zátěž systému; bez mzdové konvergence tak narůstá strukturální deficitnost průběžného financování.
Posílení kolektivního vyjednávání na odvětvové úrovni, které by comprimovalo mzdové rozdělení zdola, by tedy zároveň snižovalo redistribuční nároky na důchodový systém a~zvyšovalo jeho dlouhodobou udržitelnost.

\input{texparts/python/cz_pension_solidarity}



